\chapter{PARTE INTRODUCTORIA}

\section{GENERALIDADES}


En la actualidad, el sector bancario atraviesa una profunda transformación
impulsada por la digitalización de los servicios financieros y las crecientes
expectativas de los usuarios. Este fenómeno, denominado transformación
digital, ha redefinido la manera en que las instituciones bancarias interactúan
con sus clientes, gestionan sus operaciones y desarrollan nuevos productos y
servicios. En este contexto de evolución acelerada, la arquitectura tecnológica
que sustenta las interfaces digitales de usuario ha adquirido una importancia
estratégica sin precedentes, convirtiéndose en un factor diferenciador clave
para la competitividad y sostenibilidad de las entidades financieras.

La selección de tecnologías de interfaz de usuario representa una decisión
crítica con implicaciones de largo alcance para las instituciones bancarias.
Frameworks como Angular, React y Vue ofrecen diferentes paradigmas,
arquitecturas y ecosistemas que determinan aspectos técnicos como
rendimiento, seguridad y escalabilidad. En un sector altamente regulado como
el bancario, donde la confianza del cliente y la seguridad de las transacciones
son imperativos categóricos, la elección tecnológica adquiere además
dimensiones regulatorias ineludibles relacionadas con el cumplimiento
normativo, la protección de datos sensibles y la continuidad operativa. Por
ello, la presente investigación centra su modelo de decisión en dos ejes:
los \textbf{criterios técnicos} (propiedades intrínsecas del framework) y los
\textbf{criterios regulatorios} (exigencias normativas del sector financiero),
por ser los que efectivamente condicionan la viabilidad de una tecnología de
interfaz de usuario en un entorno bancario regulado.

Un factor adicional obliga a revisar el peso relativo de estos criterios: la
consolidación de herramientas de generación de código asistida por
Inteligencia Artificial (IA) --GitHub Copilot, Cursor, Claude Code y asistentes
equivalentes basados en modelos de lenguaje de gran escala-- que, según la
encuesta de Stack Overflow (2024), ya son utilizadas o planificadas por más
de tres cuartas partes de los desarrolladores profesionales encuestados. Este
fenómeno reduce la brecha de productividad y curva de aprendizaje que
históricamente diferenciaba a los frameworks por su ecosistema o la
disponibilidad de talento especializado (Peng, Kalliamvakou, Cihon, y
Demirer, 2023), desplazando el eje decisivo de la selección tecnológica desde
factores organizacionales --que la IA generativa tiende a nivelar-- hacia las
propiedades técnicas intrínsecas del framework y su compatibilidad con las
exigencias regulatorias del sector, que la IA no puede sustituir ni flexibilizar.
Esta reconfiguración del problema de decisión constituye el eje central que
justifica el alcance de la presente tesis.


\section{DESCRIPCIÓN DE LA PROBLEMÁTICA}


\subsection*{Situación actual en instituciones bancarias}

Un estudio reciente realizado por Deloitte (2023) entre 87 instituciones
financieras globales reveló que el 78\% de las entidades bancarias considera
crítica la modernización de sus interfaces digitales para mantener su
competitividad. Sin embargo, el mismo estudio identifica que solo el 31\%
reporta tener procesos estructurados para evaluar y seleccionar tecnologías
de interfaz de usuario, evidenciando una brecha significativa entre
necesidades estratégicas y capacidades metodológicas.

El análisis de 35 casos de implementación de plataformas digitales en bancos
latinoamericanos durante el período 2020-2024 muestra que:

\begin{itemize}
  \item 47\% excedió su presupuesto inicial en más del 40\%
  \item 62\% experimentó retrasos significativos en el despliegue (>6 meses)
  \item 28\% requirió una refactorización mayor dentro de los primeros 18 meses
  \item 53\% reportó dificultades para integrar sistemas legacy con las nuevas
        tecnologías de interfaz de usuario
\end{itemize}

Estos indicadores apuntan a deficiencias sistemáticas en los procesos de
selección tecnológica que generan impactos financieros y operativos
cuantificables.

\subsection*{Especificidad de requisitos no atendida}

Las instituciones bancarias operan bajo un conjunto de restricciones y
requerimientos particulares --predominantemente técnicos y regulatorios--
que no son adecuadamente contemplados por metodologías genéricas de
selección tecnológica:

\begin{enumerate}
  \item \textbf{Requisitos regulatorios estrictos:} el cumplimiento de normativas
        como PCI-DSS, GDPR, Basilea III y regulaciones locales impone
        restricciones específicas sobre las tecnologías de interfaz de usuario
        que deben evaluarse sistemáticamente. Un análisis del Banco
        Interamericano de Desarrollo (2024) identificó que las sanciones por
        incumplimiento normativo relacionado con interfaces digitales
        bancarias aumentaron un 215\% entre 2020-2024.
  \item \textbf{Patrones de tráfico altamente variables:} los sistemas bancarios
        experimentan fluctuaciones de carga específicas (cierres de mes, días de
        pago, campañas promocionales) que demandan características
        técnicas particulares de escalabilidad y rendimiento. Datos compilados
        por McKinsey (2023) muestran que los picos de tráfico en plataformas
        bancarias pueden superar hasta 20 veces el volumen normal, comparado
        con 5-7 veces en otros sectores.
  \item \textbf{Requisitos de integración complejos:} las arquitecturas bancarias
        típicamente incluyen entre 12-18 sistemas core interconectados con
        antigüedad promedio de 15+ años (Gartner, 2024), generando desafíos
        técnicos de integración que no se presentan en otros sectores y que la
        normativa exige documentar y auditar.
  \item \textbf{Criticidad operacional:} la indisponibilidad de plataformas
        bancarias genera impactos económicos y reputacionales excepcionales,
        además de infracciones a los acuerdos regulatorios de continuidad de
        servicio. Un estudio de Forrester (2023) cuantificó que cada minuto de
        caída en interfaces digitales bancarias representa pérdidas promedio de
        US\$67,000, significativamente superior a otros sectores.
\end{enumerate}

\subsection*{Evolución acelerada y obsolescencia tecnológica}

El ecosistema de tecnologías de interfaz de usuario experimenta una
evolución particularmente acelerada que presenta desafíos técnicos
específicos para instituciones bancarias.

\begin{table}[htbp]
\centering
\caption{Comparación de frameworks}
\label{tab:1.1}
\begin{tabularx}{\textwidth}{@{}lXXX@{}}
\toprule
Framework & Versiones Mayores (2020-2024) & Cambios Estructurales & Actualizaciones de Seguridad \\
\midrule
Angular & 6 (v9-v14)  & 3 & 27 \\
React   & 3 (v16-v18) & 2 & 21 \\
Vue     & 2 (v2-v3)   & 2 & 18 \\
\bottomrule
\end{tabularx}
\end{table}

Este ritmo evolutivo contrasta con los ciclos de implementación en
instituciones bancarias, donde el tiempo promedio entre renovaciones
tecnológicas de interfaz de usuario es de 4.7 años (FinTech Magazine, 2024),
creando una tensión estructural entre evolución tecnológica y estabilidad
operativa que un modelo de decisión técnico-regulatorio debe absorber.

\subsection*{Fragmentación decisional en criterios técnicos y regulatorios}

Un análisis de procesos de selección tecnológica en 15 instituciones bancarias
(PwC, 2023) identificó patrones problemáticos recurrentes:

\begin{itemize}
  \item \textbf{Inconsistencia metodológica:} 87\% utilizó criterios técnicos
        diferentes entre departamentos para seleccionar tecnologías de
        interfaz de usuario, sin una matriz de ponderación común.
  \item \textbf{Documentación insuficiente:} 73\% carecía de documentación
        estructurada que justificara la decisión ante auditorías regulatorias.
  \item \textbf{Evaluación parcial:} 62\% evaluó solo aspectos técnicos de
        rendimiento sin verificar sistemáticamente el cumplimiento normativo.
  \item \textbf{Ausencia de métricas:} 91\% no estableció KPIs técnicos ni
        regulatorios específicos para medir el éxito de la selección tecnológica.
  \item \textbf{Desconexión histórica:} 78\% no incorporó sistemáticamente
        lecciones aprendidas de implementaciones anteriores.
\end{itemize}

Esta fragmentación obstaculiza el desarrollo de capacidades institucionales
sostenibles y reproduce ineficiencias entre proyectos.


\section{FORMULACIÓN DEL PROBLEMA GENERAL}


¿De qué manera puede desarrollarse una metodología basada en criterios
técnicos y regulatorios para la selección de tecnologías de interfaz de usuario
(Angular, React y Vue) en instituciones bancarias, que sea coherente con el
contexto actual de generación de código asistida por Inteligencia Artificial, y
cuya aplicación sistemática resulte en decisiones tecnológicas mejor
alineadas con las necesidades específicas del sector?


\subsection{Formulación de los problemas específicos}

\subsubsection{Problema específico 1}

¿Qué criterios técnicos y regulatorios específicos del sector bancario
        deben ser incorporados en la evaluación comparativa de tecnologías de
        interfaz de usuario para maximizar su alineación con requisitos
        sectoriales?

\subsubsection{Problema específico 2}

¿Qué instrumentos metodológicos permiten cuantificar objetivamente la
        idoneidad de diferentes frameworks de interfaz de usuario (Angular,
        React y Vue) frente a los requisitos técnicos y regulatorios de
        instituciones bancarias?

\subsubsection{Problema específico 3}

¿Qué proceso estructurado de toma de decisiones permite integrar
        coherentemente las dimensiones técnica y regulatoria en la selección
        de tecnologías de interfaz de usuario para instituciones financieras
        con arquitecturas heterogéneas y requisitos regulatorios estrictos?

\subsubsection{Problema específico 4}

¿Cómo puede validarse empíricamente la efectividad de la metodología
        propuesta en escenarios reales de transformación digital bancaria?

\subsubsection{Problema específico 5}

¿Qué indicadores permiten medir objetivamente el éxito de una
        selección tecnológica de interfaz de usuario en el contexto específico
        de implementaciones bancarias?

\subsubsection{Problema específico 6}

¿De qué manera la disponibilidad de herramientas de generación de
        código asistida por Inteligencia Artificial modifica la relevancia
        relativa de los criterios técnicos y regulatorios frente a los
        criterios organizacionales tradicionalmente considerados en la
        selección de tecnologías de interfaz de usuario?

\section{JUSTIFICACIÓN DEL PROBLEMA}

\subsection*{Impacto económico y competitivo}

La problemática descrita genera impactos cuantificables:

\begin{itemize}
  \item \textbf{Costos directos:} según Boston Consulting Group (2023), las
        decisiones subóptimas en selección de tecnologías de interfaz de
        usuario representan un sobrecosto promedio del 32\% en proyectos de
        transformación digital bancaria.
  \item \textbf{Time-to-market:} las entidades con procesos maduros de
        selección tecnológica reducen en 47\% el tiempo de despliegue de
        nuevas funcionalidades (EY, 2024).
  \item \textbf{Riesgo regulatorio:} las instituciones sin un proceso técnico y
        regulatorio documentado experimentan hallazgos de auditoría
        relacionados con interfaces digitales con una frecuencia 2.3 veces
        superior al promedio sectorial (Stack Overflow Banking Survey, 2023).
  \item \textbf{Ventaja competitiva:} las instituciones con metodologías
        estructuradas de selección tecnológica presentan un Net Promoter Score
        promedio 12 puntos superior en canales digitales (J.D. Power Banking
        Satisfaction Study, 2024).
\end{itemize}

\subsection*{El efecto nivelador de la IA generativa sobre los factores organizacionales}

Hasta hace pocos años, criterios organizacionales como la disponibilidad de
talento especializado o la curva de aprendizaje de un framework pesaban de
forma comparable a los criterios técnicos en las decisiones tecnológicas
(Mishra y Mishra, 2021). La evidencia reciente sobre productividad de
desarrolladores con asistentes de código basados en IA --un incremento de
hasta 55\% en la velocidad de completado de tareas de programación,
documentado por Peng et al. (2023) en un experimento controlado con GitHub
Copilot-- sugiere que esta brecha se está reduciendo: un desarrollador con
experiencia media, asistido por IA, puede volverse productivo en un framework
que no domina en un tiempo significativamente menor al histórico. Esto no
elimina la necesidad de criterios de evaluación, pero sí modifica cuáles son
los verdaderamente determinantes: como se detalla en el Capítulo 2 (sección
2.11.3) y se valida empíricamente en el Capítulo 4, el código generado por IA
es más consistente y auditable en frameworks fuertemente tipados y
opinionados, lo que convierte la \emph{calidad y verificabilidad del código
asistido por IA} en un criterio técnico de primer orden, y la
\emph{gobernanza regulatoria del uso de herramientas de IA generativa} en un
criterio regulatorio emergente --ninguno de los dos contemplado por las
metodologías genéricas de evaluación tecnológica existentes (Jansen y Lee,
2011).


\section{OBJETIVOS DEL ESTUDIO}

\subsection{Objetivo general}


Desarrollar y validar una metodología basada en criterios técnicos y
regulatorios para la selección de tecnologías de interfaz de usuario (Angular,
React y Vue) en el contexto de transformación digital de instituciones
bancarias, incorporando el efecto de las herramientas de generación de
código asistida por Inteligencia Artificial en la relevancia relativa de dichos
criterios, y optimizando la alineación entre elecciones tecnológicas y
necesidades institucionales.


\subsection{Objetivos específicos}

\subsubsection{Objetivo específico 1}

Identificar y categorizar los criterios técnicos y regulatorios críticos
        del sector bancario que deben considerarse en la evaluación de
        tecnologías de interfaz de usuario, mediante revisión sistemática de
        literatura especializada y consulta a expertos sectoriales.

\subsubsection{Objetivo específico 2}

Diseñar instrumentos metodológicos cuantitativos y cualitativos para la
        evaluación comparativa objetiva de Angular, React y Vue frente a
        requisitos técnicos y regulatorios bancarios, incluyendo matrices de
        decisión ponderadas, rúbricas de evaluación y mecanismos de
        benchmarking técnico.

\subsubsection{Objetivo específico 3}

Formular un proceso estructurado de toma de decisiones que integre
        coherentemente las dimensiones técnica y regulatoria en la selección
        de tecnologías de interfaz de usuario para instituciones financieras,
        con definición de roles, responsabilidades y flujos de trabajo
        específicos para arquitecturas heterogéneas y entornos regulatorios
        estrictos.

\subsubsection{Objetivo específico 4}

Implementar y validar la metodología propuesta mediante su aplicación
        en un caso de estudio real en una institución bancaria en proceso de
        transformación digital, documentando sistemáticamente resultados,
        ajustes metodológicos y lecciones aprendidas.

\subsubsection{Objetivo específico 5}

Desarrollar un sistema de indicadores técnicos y regulatorios para
        evaluar el éxito de la selección tecnológica de interfaz de usuario en
        implementaciones bancarias.

\subsubsection{Objetivo específico 6}

Determinar el efecto de la adopción de herramientas de generación de
        código asistida por Inteligencia Artificial sobre la ponderación relativa
        de los criterios técnicos y regulatorios frente a los criterios
        organizacionales, mediante análisis de sensibilidad comparado con y
        sin dichos criterios en el modelo.

\section{HIPÓTESIS DEL ESTUDIO}

\subsection{Hipótesis general}


La implementación de una metodología estructurada basada en criterios
técnicos y regulatorios --coherente con el contexto de generación de código
asistida por Inteligencia Artificial-- para la selección de tecnologías de
interfaz de usuario en instituciones bancarias, resulta en decisiones
tecnológicas significativamente más alineadas con los requisitos particulares
de las entidades financieras y genera mejores resultados en términos de
rendimiento técnico y cumplimiento normativo, en comparación con enfoques
de selección tecnológica genéricos, no estructurados o centrados
predominantemente en factores organizacionales.


\subsection{Hipótesis específicas}

\subsubsection{Hipótesis específica 1}

La identificación y ponderación sistemática de criterios técnicos y
        regulatorios específicos del sector bancario para la evaluación de
        tecnologías de interfaz de usuario permite una valoración más precisa
        de la idoneidad de frameworks como Angular, React y Vue frente a los
        requerimientos particulares de instituciones financieras, en
        comparación con la aplicación de criterios generales de evaluación
        tecnológica.

        \textbf{Variables:}
        \begin{itemize}
          \item \textit{Variable independiente:} Aplicación de criterios
                técnicos y regulatorios específicos del sector bancario
                ponderados sistemáticamente.
          \item \textit{Variable dependiente:} Precisión en la valoración de
                idoneidad de frameworks de interfaz de usuario para requisitos
                financieros.
          \item \textit{Indicadores:} Grado de alineación entre características
                tecnológicas y requisitos técnicos/regulatorios bancarios; nivel
                de satisfacción de stakeholders con el proceso evaluativo.
        \end{itemize}


\subsubsection{Hipótesis específica 2}

La utilización de instrumentos metodológicos cuantitativos y
        cualitativos diseñados específicamente para evaluar tecnologías de
        interfaz de usuario en contextos bancarios incrementa
        significativamente la objetividad, reproducibilidad y transparencia del
        proceso de selección tecnológica, comparado con evaluaciones basadas
        en criterios subjetivos o instrumentos genéricos.

        \textbf{Variables:}
        \begin{itemize}
          \item \textit{Variable independiente:} Utilización de instrumentos
                metodológicos especializados (cuantitativos y cualitativos).
          \item \textit{Variable dependiente:} Objetividad, reproducibilidad y
                transparencia del proceso de selección.
          \item \textit{Indicadores:} Consistencia de resultados entre
                evaluadores independientes; trazabilidad de decisiones;
                variabilidad de calificaciones para criterios equivalentes.
        \end{itemize}


\subsubsection{Hipótesis específica 3}

La implementación de un proceso decisional estructurado que integre
        exclusivamente consideraciones técnicas y regulatorias, con roles y
        responsabilidades claramente definidos, optimiza la selección de
        tecnologías de interfaz de usuario para arquitecturas bancarias
        heterogéneas, resultando en mejores decisiones organizacionales
        comparado con aproximaciones decisionales fragmentadas o
        excesivamente centradas en preferencias subjetivas de los equipos.

        \textbf{Variables:}
        \begin{itemize}
          \item \textit{Variable independiente:} Implementación de un proceso
                decisional estructurado centrado en criterios técnicos y
                regulatorios, con roles definidos.
          \item \textit{Variable dependiente:} Calidad de integración entre
                consideraciones técnicas y regulatorias.
          \item \textit{Indicadores:} Tiempo requerido para toma de decisiones;
                nivel de participación cross-funcional; grado de consenso
                organizacional; documentación de justificaciones decisionales.
        \end{itemize}


\subsubsection{Hipótesis específica 4}

La validación empírica de la metodología propuesta mediante su
        aplicación en un caso de estudio real en una institución bancaria
        demuestra su efectividad práctica y aplicabilidad, generando resultados
        cualitativa y cuantitativamente mejores que los obtenidos con
        metodologías previas no especializadas.

        \textbf{Variables:}
        \begin{itemize}
          \item \textit{Variable independiente:} Aplicación de la metodología
                propuesta en caso real.
          \item \textit{Variable dependiente:} Efectividad práctica y
                aplicabilidad de la metodología.
          \item \textit{Indicadores:} Comparativa entre resultados pre/post
                implementación; nivel de adherencia metodológica en
                implementación real; adaptaciones requeridas en
                implementación.
        \end{itemize}


\subsubsection{Hipótesis específica 5}

El sistema de indicadores técnicos y regulatorios desarrollado para
        evaluar el éxito de selecciones tecnológicas de interfaz de usuario en
        implementaciones bancarias permite medir con mayor precisión el
        impacto técnico, operativo y de cumplimiento normativo de las
        decisiones tecnológicas, en comparación con métricas genéricas de
        evaluación de proyectos.

        \textbf{Variables:}
        \begin{itemize}
          \item \textit{Variable independiente:} Sistema de indicadores
                técnicos y regulatorios para contexto bancario.
          \item \textit{Variable dependiente:} Precisión en la medición de
                impactos técnicos y de cumplimiento.
          \item \textit{Indicadores:} Sensibilidad y especificidad de métricas;
                correlación entre indicadores propuestos y resultados
                organizacionales; capacidad predictiva de los indicadores sobre
                éxito implementativo.
        \end{itemize}


\subsubsection{Hipótesis específica 6}

La disponibilidad de herramientas de generación de código asistida
        por Inteligencia Artificial reduce significativamente la relevancia
        relativa de los criterios organizacionales (disponibilidad de talento,
        ecosistema, curva de aprendizaje) frente a los criterios técnicos y
        regulatorios en la selección de tecnologías de interfaz de usuario en
        instituciones bancarias, sin alterar --e incluso reforzando-- la
        preferencia por frameworks con arquitecturas tipadas y opinionadas.

        \textbf{Variables:}
        \begin{itemize}
          \item \textit{Variable independiente:} Grado de adopción de
                herramientas de generación de código asistida por IA en el
                equipo de desarrollo.
          \item \textit{Variable dependiente:} Relevancia relativa de los
                criterios técnicos y regulatorios frente a los organizacionales
                en la puntuación final de la decisión.
          \item \textit{Indicadores:} Variación de la brecha de puntuación
                ($\Delta S$) entre frameworks al incluir o excluir los criterios
                asociados a la IA generativa; proporción del peso Delphi
                asignado a criterios técnicos y regulatorios frente al
                históricamente asignado a criterios organizacionales.
        \end{itemize}

\section{METODOLOGÍA}

\subsection{Tipo y nivel de investigación}

La metodología seguida para el desarrollo de la presente tesis es la siguiente:

\textbf{Enfoque de investigación.} La investigación requiere un enfoque
        mixto (cuantitativo-cualitativo) dado que se necesita identificar y
        categorizar criterios técnicos y regulatorios específicos del sector
        bancario, diseñar instrumentos metodológicos para evaluación objetiva,
        determinar el efecto moderador de la IA generativa sobre dichos
        criterios, y validar empíricamente la efectividad de la metodología
        propuesta.



\textbf{Diseño de investigación.} Se propone un diseño de investigación
        aplicada con componente exploratorio y descriptivo seguido de una fase
        experimental a través de caso de estudio:
        \begin{itemize}
          \item Fase exploratoria-descriptiva: para identificar criterios
                técnicos y regulatorios relevantes en la selección de
                tecnologías de interfaz de usuario en el sector bancario, y
                para caracterizar el efecto de las herramientas de IA
                generativa sobre la relevancia de criterios organizacionales.
          \item Fase experimental: para validar la metodología propuesta
                mediante un caso de estudio real en una institución bancaria,
                incluyendo un análisis de sensibilidad con y sin los criterios
                asociados a la IA generativa.
        \end{itemize}


\subsection{Variables de estudio}

Las variables independientes y dependientes de cada hipótesis, junto con sus
indicadores y los instrumentos de medición asociados, se presentan integrados
por hipótesis en la matriz de consistencia (sección 1.8), de modo que la
correspondencia problema--objetivo--hipótesis--variable--indicador--método
quede trazada en un solo instrumento.


\subsection{Técnicas e instrumentos de recolección de datos}

\textbf{Métodos de recolección de datos.}
        \begin{itemize}
          \item \emph{Para identificación de criterios (PE1):} revisión
                sistemática de literatura, entrevistas semiestructuradas con
                expertos del sector y técnica Delphi para consensuar y validar
                los criterios técnicos y regulatorios críticos con panel de
                expertos.
          \item \emph{Para diseño de instrumentos metodológicos (PE2):}
                construcción de matrices de evaluación, escalas ponderadas y
                modelos de puntuación cuantitativa.
          \item \emph{Para el proceso decisional técnico-regulatorio (PE3):}
                mapeo de procesos, estudios de caso comparativos y
                entrevistas a stakeholders.
          \item \emph{Para validación empírica (PE4):} caso de estudio,
                experimentación controlada y medición pre/post-implementación.
          \item \emph{Para indicadores de éxito (PE5):} establecimiento de KPIs
                técnicos y regulatorios, tableros de medición y análisis de
                impacto.
          \item \emph{Para el efecto de la IA generativa (PE6):} análisis de
                sensibilidad comparado (modelo con y sin criterios de IA) y
                encuesta complementaria al panel Delphi sobre su percepción
                del impacto de los asistentes de código en la relevancia de
                criterios organizacionales.
        \end{itemize}



\textbf{Instrumentos de investigación.} Cuestionarios estructurados
        para expertos, guías de entrevista semiestructurada, matrices de
        evaluación comparativa, plantillas de documentación de procesos,
        herramientas de análisis cuantitativo y cualitativo, y frameworks de
        medición de indicadores.


\subsection{Población y muestra}

\textbf{Población y muestra.} Población: instituciones bancarias con
        necesidades de selección de tecnologías de interfaz de usuario.
        Muestra: selección estratificada de instituciones de diferentes
        tamaños y tipos (banca comercial, inversión, etc.). Muestreo no
        probabilístico por conveniencia y juicio de expertos.



\textbf{Análisis de datos.} Análisis cualitativo mediante codificación
        temática para entrevistas y feedback; análisis cuantitativo mediante
        estadística descriptiva e inferencial para validación de la efectividad;
        triangulación metodológica para asegurar la validez de los hallazgos.


\section{MATRIZ DE CONSISTENCIA}

\begin{landscape}
\newgeometry{left=2cm,right=2cm,top=2cm,bottom=2cm}
\begin{spacing}{1}
\footnotesize
\begin{longtable}{@{}p{3.9cm}p{3.9cm}p{3.9cm}p{3.95cm}p{3.95cm}p{3.95cm}@{}}
\caption{Matriz de Consistencia}\label{tab:1.2}\\
\toprule
Problemas & Objetivos & Hipótesis & Variables & Indicadores & Metodología \\
\midrule
\endfirsthead
\multicolumn{6}{l}{\small\itshape (Tabla \ref{tab:1.2} -- continuación)}\\
\toprule
Problemas & Objetivos & Hipótesis & Variables & Indicadores & Metodología \\
\midrule
\endhead
\midrule
\multicolumn{6}{r}{\small\itshape Continúa en la página siguiente}\\
\endfoot
\bottomrule
\endlastfoot
PG: ¿De qué manera puede desarrollarse una metodología basada en criterios
técnicos y regulatorios para la selección de tecnologías de interfaz de
usuario en instituciones bancarias, coherente con el contexto de IA
generativa?
&
OG: Desarrollar y validar una metodología basada en criterios técnicos y
regulatorios para la selección de tecnologías de interfaz de usuario en
instituciones bancarias, incorporando el efecto de la IA generativa.
&
HG: La implementación de una metodología estructurada basada en criterios
técnicos y regulatorios, coherente con el contexto de IA generativa, resulta
en decisiones tecnológicas significativamente más alineadas con los
requisitos particulares de las entidades financieras, en comparación con
enfoques genéricos o centrados predominantemente en factores
organizacionales.
&
VI: Metodología técnico-regulatoria de selección \newline
VD: Calidad de la decisión tecnológica
&
Alineación con requisitos técnicos y normativos \newline
Robustez de la decisión ($\Delta S$)
&
Enfoque mixto con triangulación metodológica \newline
Delphi + AHP \newline
Caso de estudio \\
\addlinespace
PE1: ¿Qué criterios técnicos y regulatorios específicos del sector bancario
deben ser incorporados en la evaluación?
&
OE1: Identificar y categorizar los criterios técnicos y regulatorios críticos
del sector bancario para la evaluación de tecnologías de interfaz de usuario.
&
HE1: La identificación y ponderación sistemática de criterios técnicos y
regulatorios permite una valoración más precisa de la idoneidad de
frameworks.
&
VI: Aplicación de criterios técnicos y regulatorios bancarios \newline
VD: Precisión en valoración de idoneidad
&
Grado de alineación entre características y requisitos \newline
Nivel de satisfacción de stakeholders
&
Matriz de trazabilidad requisitos-características \newline
Encuestas de satisfacción a stakeholders \\
\addlinespace
PE2: ¿Qué instrumentos metodológicos permiten cuantificar objetivamente la
idoneidad de diferentes frameworks?
&
OE2: Diseñar instrumentos metodológicos cuantitativos y cualitativos para
evaluación comparativa objetiva.
&
HE2: La utilización de instrumentos metodológicos especializados incrementa
significativamente la objetividad, reproducibilidad y transparencia.
&
VI: Utilización de instrumentos metodológicos especializados \newline
VD: Objetividad, reproducibilidad y transparencia
&
Consistencia entre evaluadores \newline
Trazabilidad de decisiones
&
Test de concordancia inter-evaluadores \newline
Auditoría de trazabilidad decisional \\
\addlinespace
PE3: ¿Qué proceso estructurado de toma de decisiones integra
coherentemente las dimensiones técnica y regulatoria?
&
OE3: Formular un proceso estructurado de toma de decisiones que integre
las dimensiones técnica y regulatoria.
&
HE3: La implementación de un proceso decisional estructurado
técnico-regulatorio optimiza la integración de consideraciones frente a
aproximaciones fragmentadas.
&
VI: Proceso decisional técnico-regulatorio estructurado \newline
VD: Calidad de integración de consideraciones
&
Tiempo para decisiones \newline
Participación cross-funcional
&
Registro de tiempos decisionales \newline
Análisis de participación por área \\
\addlinespace
PE4: ¿Cómo puede validarse empíricamente la efectividad de la metodología
propuesta?
&
OE4: Implementar y validar la metodología propuesta mediante caso de
estudio real.
&
HE4: La validación empírica mediante caso de estudio demuestra efectividad
práctica y aplicabilidad.
&
VI: Aplicación de metodología en caso real \newline
VD: Efectividad práctica y aplicabilidad
&
Comparativa pre/post implementación \newline
Adherencia metodológica
&
Análisis comparativo de resultados \newline
Lista de verificación de cumplimiento metodológico \\
\addlinespace
PE5: ¿Qué indicadores permiten medir objetivamente el éxito de una
selección tecnológica de interfaz de usuario?
&
OE5: Desarrollar un sistema de indicadores técnicos y regulatorios para
evaluar el éxito de la selección tecnológica.
&
HE5: El sistema de indicadores técnicos y regulatorios permite medir con
mayor precisión el impacto técnico y de cumplimiento normativo.
&
VI: Sistema de indicadores técnicos y regulatorios \newline
VD: Precisión en medición de impactos
&
Sensibilidad y especificidad de métricas \newline
Correlación indicadores-resultados
&
Análisis de sensibilidad y especificidad \newline
Estudios de correlación estadística \\
\addlinespace
PE6: ¿De qué manera la disponibilidad de herramientas de IA generativa
modifica la relevancia relativa de los criterios técnicos y regulatorios frente
a los organizacionales?
&
OE6: Determinar el efecto de la adopción de herramientas de IA generativa
sobre la ponderación relativa de los criterios técnicos y regulatorios
frente a los organizacionales.
&
HE6: La disponibilidad de herramientas de IA generativa reduce la
relevancia relativa de los criterios organizacionales frente a los técnicos y
regulatorios, sin alterar la preferencia por frameworks tipados y
opinionados.
&
VI: Adopción de herramientas de IA generativa en desarrollo \newline
VD: Relevancia relativa de criterios técnicos-regulatorios
&
Variación de $\Delta S$ con/sin criterios de IA \newline
Proporción del peso Delphi técnico-regulatorio
&
Análisis de sensibilidad comparado (sección 4.3) \newline
Encuesta al panel Delphi sobre percepción del impacto de IA \\
\end{longtable}
\normalsize
\end{spacing}
\restoregeometry
\end{landscape}

